\begin{description}
	\item[1行目] 読み込み他，ファイル操作を便利にするパッケージ\texttt{tidyverse}を読み込みます\footnote{\texttt{tidyverse}パッケージがない人はインストールしてください。関連パッケージをいろいろまとめたパッケージ群パッケージなので，少し時間がかかります。}。
	\item[2-9行目] データファイルを変形し，\texttt{dataset}オブジェクトに代入しています。ここで\verb| %>% | はパイプ演算子と呼ばれるもので，左の操作を右の操作へつなげる役割をします。
		\begin{description}
			\item[3行目] データを2020年度のものだけに絞るため，\verb|dplyr|パッケージの\verb|filter|関数を適用しています。
			\item[4行目] データを投手のものだけに絞るため，\verb|dplyr|パッケージの\verb|filter|関数を適用しています。
			\item[5行目] 変数を必要なものだけに絞るため，\verb|dplyr|パッケージの\verb|select|関数を適用しています。
			\item[6行目] データセットから欠損値を除いています。
			\item[7行目] 表示用に，データを勝利数(変数\verb|Win|)の多い順に並べ直しています。
			\item[8行目] 以下の計算に使う変数だけに絞り込んでいます。
			\item[9行目] データセットを行列方に変換しています。
		\end{description}
